
\documentclass[12pt]{article}

% --- Página y tipografía ---
\usepackage[letterpaper,margin=2.5cm]{geometry}
\usepackage[T1]{fontenc}
\usepackage[utf8]{inputenc} % si compilas con pdfLaTeX
\usepackage{lmodern}
\usepackage{microtype}

% --- Imágenes y color ---
\usepackage{graphicx}
\usepackage{xcolor}

% --- Control fino de espacios ---
\usepackage{setspace}
\setlength{\parindent}{0pt}

\begin{document}
\thispagestyle{empty}

% ===== Encabezado con logos + texto =====
\begin{minipage}[c]{0.18\textwidth}
    \centering
    % Cambia por tu logo izquierdo
    \includegraphics[width=0.95\linewidth]{logo_usac.jpeg}
\end{minipage}
\hfill
\begin{minipage}[c]{0.60\textwidth}
    \small
    Universidad de San Carlos de Guatemala\\
    Escuela de Ciencias Físicas y Matemáticas\\
    Nombre estudiante\\
    Carnet: \\
    Programación 1\\
\end{minipage}
\hfill
\begin{minipage}[c]{0.18\textwidth}
    \centering
    % Cambia por tu logo derecho
    \includegraphics[width=1.4\linewidth]{logo_ecfm.jpg}
\end{minipage}

\vspace{0.5cm}

% Línea horizontal superior (gruesa)
\noindent\rule{\textwidth}{1.2pt}

\vspace{0.2cm}

% ===== Título =====
\begin{center}
    {\Large\scshape Ensayo}\\[0.3em]
\end{center}

\vspace{0.1cm}

% Fecha
\begin{center}
    \small\scshape viernes 06 de febrero de 2026
\end{center}

\vspace{0.2cm}

% Línea horizontal inferior (gruesa)
\noindent\rule{\textwidth}{1.2pt}

\vspace{0.6cm}

% ===== Caja de resumen =====
\noindent
\colorbox{gray!36}{%
    \parbox{\textwidth}{%
        \vspace{0.6em}
        \textbf{Resumen}\\[0.4em]
        \small
El presente ensayo se enfoca en  dara respuesta a las siguientes preguntas:  ¿cuál sería tú área de investigación en Física? y ¿explica por qué la programación te sería útil en dicho campo?, se realizá una descripción para cada pregunta, desde el punto de vista de un estudiante de física. 
        \vspace{0.8em}
    }%
}


\section{Pregunta No. 1}
\subsection{¿Cuál sería tú área de investigación en Física?}

Es un tema facinante, si fuera un investigador  de física, el área de investigación sería la Física de Radiociones, ya que es un área que conbina las leyes naturales y la vida y poder contribuir a savar las mismas.\\
 Y en esta misma área el interés en la investigación sería en Radioterapia que es el uso de radiación ionizante para tratar tumores (física de partículas, dosis que se requiere a un paciente).\\
 De acuerdo al siguiente texto "La radiación ionizante es la liberación de electrones mediante la emisión de energía en forma de ondas y partículas que se producen de forma natural en los materiales que conocemos como radiactivos (suelo, agua y vegetación) y de forma artificial en isótopos creados por el hombre, tal como los equipos que producen Rayos-X" (ATSDR en Español, 2016).\\
Por tanto comprender la liberación de los electrones puede alterar o modificar la estructura molecular de un ser vivo.


\section{Pregunta 2}
\subsection{¿Explica por qué la programación te sería útil en dicho campo?}

En física médica, no siempre se puedes experimentar directamente con humanos por razones éticas y de seguridad.\\

¿Para qué sirve? Se usan algoritmos (como Geant4 o MCNP) para simular millones de trayectorias de partículas (fotones, electrones, protones) golpeando un tejido.\\

La programación te permite: Definir la geometría del paciente y calcular exactamente cuánta energía se deposita en cada milímetro cúbico.



\end{document}

\end{document}